%
%
\section{Stochastic Matrices over PCA-enriched categories}\label{sec:cmatrix}    
It is well-known (see, e.g., \cite[Exercises VIII.2.5-6]{mac_lane_categories_1978}) that, from a category enriched over commutative monoids, one can freely generate the finite biproduct category of its matrices. In this section, we show a similar construction for $\Cat{PCA}$-enriched categories: given such a category $\Cat{C}$, one can construct the convex biproduct category $\stmat{\Cat{C}}$ of \emph{stochastic matrices} over $\Cat{C}$.

Objects of $\stmat{\Cat{C}}$ are words in $Ob(\Cat{C})^*$. We will write $\bigoplus_{k=1}^m U_k$ for the word $U_1\dots U_m$ and $\zero$ for the empty word. We will denote objects of $\Cat{C}$ by $U,V$ and those of $\stmat{\Cat{C}}$ by $P,Q$. In $\stmat{\Cat{C}}$, an arrow $M\colon\bigoplus_{k=1}^n U_k\to \bigoplus_{k=1}^m V_k$ is an equivalence class of  $m\times n$ matrices with $(j,i)$-entries given by pairs $ (p_{ji}, f_{ji})$ where $f_{ji}\in\Cat{C}[U_i,V_j]$ and $p_{ji}\in[0,1]$ satisfy $\sum_{j=1}^{m}p_{ji}\le 1$. Two matrices $M$ and $M'$ are equivalent, in symbols $M\equiv M'$, if $p_{ji}\cdot f_{ji}= p'_{ji}\cdot f'_{ji}$ in $\Cat{C}[U_i,V_j]$ for all $i,j$. 
\begin{remark}
The use of pairs $ (p_{ji}, f_{ji})$ as entries of matrices is necessary to specify the constraints $\sum_{j=1}^{m}p_{ji}\le 1$. However, $\equiv$ ensures that one can safely write $M_{ji}= p_{ji}\cdot f_{ji}$.
\end{remark}

The composition of two morphisms $M\colon\bigoplus_{k=1}^n U_k\to \bigoplus_{k=1}^m V_k$ and $M'\colon\bigoplus_{k=1}^m V_k\to \bigoplus_{k=1}^l W_k$ %, given by an $l\times m$ matrix with entries $N_{uj}=q_{uj}\cdot g_{uj}$,  
is obtained by matrix multiplication $M'M$: for all $u\in \{1,\dots,l\}$ and $i\in \{1,\dots,n\}$ the entry at $(u,i)$ is 
\begin{equation}\label{eq:matrixmult}(M'M)_{ui} \defeq  r_{ui}\cdot \left(\sum_{j=1}^{m} \frac{p'_{uj}p_{ji}}{r_{ui}}\cdot (f_{ji};f'_{uj})\right) = \sum_{j=1}^{m} p'_{uj}p_{ji}\cdot (f_{ji};f'_{uj})\end{equation} 
where the convex sums $\sum_k p_k\cdot (-)_k$ are those provided by the $\Cat{PCA}$ enrichment of $\Cat{C}$, the composition $;$ is in $\Cat{C}$ and
$r_{ui}= (\sum_{j=1}^{m} p'_{uj}p_{ji})$. Simple computations confirm that $\sum_{u=1}^{l}r_{ui}\leq 1$ and $\frac{p'_{uj}p_{ji}}{r_{ui}}\in[0,1]$. Note that the rightmost equality in \eqref{eq:matrixmult} follows from Lemma~\ref{lemma:initialproperties2}.\ref{lemma q per somma}.
%The composition is well defined since $\sum_{u=1}^{l}\sum_{j=i}^{m} q_{uj}p_{ji}=\sum_{j=i}^{m}p_{ji}\sum_{u=1}^{l}q_{uj}\le 1$ and $\frac{q_{uj}p_{ji}}{\sum_{j=i}^{m} q_{uj}p_{ji}}\in[0,1]$. 
For all objects $P=\bigoplus_{k=1}^n U_k$, $\id{P}$ is the $n\times n$ matrix with entries $(\id{P})_{jj}=1\cdot \id{U_j}$ and, for $i\neq j$, $(\id{P})_{ji}=0 \cdot \star$.

\begin{example}\label{ex:matrices}
Recall $\Cat{1}^+$ and $(A^*)^+$ from Example~\ref{ex:PCA}.
In $\stmat{\Cat{1}^+}$ objects are natural numbers and arrows $n \to m$ are the usual $m\times n$ sub-stochastic matrices (i.e., $\sum_{j=1}^mp_{ji}\leq 1$).
%
In $\stmat{(A^*)^+}$ objects are natural numbers and arrows $n \to m$ are $m\times n$ matrices with entries in $\Dis(A^*)$. Consider for instance the matrices $N\colon 2\to 3$ and $M\colon 2 \to 2$ on the left below.  The composition $M;N\colon 2 \to 3$ is the matrix on the right.
\begin{equation*}
\begin{pmatrix}
    \frac{1}{2} \cdot a & \frac{1}{2} \cdot c\\
    \frac{1}{3} \cdot ab & 0 \cdot \star \\
    0 \cdot \star & \frac{1}{3} \cdot \id{}
\end{pmatrix} 
\begin{pmatrix}
    \frac{1}{2} \cdot a & 1 \cdot c\\
    \frac{1}{2} \cdot ab & 0 \cdot \star \\
\end{pmatrix} 
=
\begin{pmatrix}
    \frac{1}{2} \cdot (\frac{1}{2} \cdot aa +\frac{1}{2}abc) & \frac{1}{2} \cdot ca\\
    \frac{1}{6} \cdot aab & \frac{1}{3} \cdot cab \\
    \frac{1}{6} \cdot ab  & 0 \cdot \star
\end{pmatrix} 
\end{equation*}
\end{example}


%
%
%$\stmat{\Cat{C}}$ has a strict symmetric monoidal structure as follows: $\bigoplus_{k=1}^n A_k \piu \bigoplus_{k=1}^{n'} A'_k$ is given by the formal $\piu$'s of the elements $A_k$ for $i=1,\dots,n$ and $A'_k$ for $i=1,\dots,n'$. For two matrices $M\colon\bigoplus_{k=1}^n A_k\to \bigoplus_{k=1}^m B_k$ and $N\colon\bigoplus_{k=1}^{n'} A'_k\to \bigoplus_{k=1}^{m'} B'_k$ is given by the $(m+m')\times (n+n')$ matrix 
%It is convenient to illustrate the underlying monoidal structure as well as the natural and coherent monoids and PCAs.
For arbitrary matrices $M$ and $N$, $M\oplus N$ is defined as below on the left
\begin{equation}\label{eq:matsmc}
M \piu N \defeq \begin{pmatrix}
    M & \emptyset\\
    \emptyset & N
\end{pmatrix}
\qquad
\symm{P}{Q}\defeq \begin{pmatrix}
    \emptyset & \id{Q}\\
    \id{P} & \emptyset
\end{pmatrix}
\end{equation}
where $\emptyset$ is the matrix with all entries $ 0\cdot \star$. The symmetry $\symmp{P}{Q}\colon P\oplus Q \to Q \oplus P$ is defined as on the right above. %Simple matrix multiplications confirms that these structures make $\stmat{\Cat{C}}$ a strict symmetric monoidal structure.
%
%$M\colon U'\to U$ and $N\colon V'\to V$ it holds $(M\piu I);\sigma_{U,V} =\sigma_{U',V};(I\piu M)$ and $(I\piu N);\sigma_{U,V}= \sigma_{U,V'};(N\piu I)$.
%
%This
%
%The unit is given by the $0\times 0$ matrix, i.e.\ the empty matrix.
%
%The symmetry for $U=\bigoplus_{k=1}^n A_k$ and $V=\bigoplus_{k=1}^m B_k$ is  given  by 
%\[\symm{U}{V}= \begin{pmatrix}
%    \emptyset_{m\times n} & I_V\\
%    I_U & \emptyset_{n\times m}
%\end{pmatrix}\]
%and a simple matrix multiplication shows that for $M\colon U'\to U$ and $N\colon V'\to V$ it holds $(M\piu I);\sigma_{U,V} =\sigma_{U',V};(I\piu M)$ and $(I\piu N);\sigma_{U,V}= \sigma_{U,V'};(N\piu I)$.
% 
%
%
Every object $P$ carries co-pca and monoid structures. These are defined for all objects $U\in \Cat{C}$ as follows, where $!_U$ (respectively $?_U$) is the unique matrix with $0$ rows (columns).
\begin{equation}\label{eq:matmonpca}
\diagp{U}\defeq\begin{pmatrix}
    p\cdot \id{U} \\
    (1-p)\cdot \id{U}
\end{pmatrix}
\qquad
\bang{U}\defeq !_U
\qquad 
\cobang{U}\defeq ?_U
\qquad
\codiag{U}\defeq \begin{pmatrix}
    1\cdot \id{U} & 1\cdot \id{U}
\end{pmatrix} 
\end{equation}
 For arbitrary objects $P$ of $\stmat{\Cat{C}}$, co-pcas and monoids are defined inductively as:
\begin{equation}\label{eq:indmoncopca}
\begin{array}{c|c}
\begin{array}{ccc}
 \diagp{\zero}\defeq \id{\zero}\; &\; \bang{\zero}\defeq \id{\zero}\; &\; \bang{U \oplus P}\defeq \bang{U}\oplus \bang{P}
 \end{array}
 &
 \begin{array}{ccc}
\codiag{\zero}\defeq \id{\zero}\; &\; \cobang{\zero}\defeq \id{\zero} \;& \;\cobang{U \oplus P}\defeq \cobang{U}\oplus \cobang{P}
\end{array} \\
\diagp{U \oplus P}\defeq (\diagp{U} \piu\, \diagp{P});(\id{U} \piu \symm{U}{P} \piu \id{P})&
\codiag{U\oplus P}\defeq (\id{U} \piu \symm{P}{U} \piu \id{P});(\codiag{U}\piu \codiag{P})
\end{array}
\end{equation}

%\[\bang{\zero}\defeq \id{\zero}\]
%\[\]
%
%
%\[\]
%\[
%
%\]
%
%\[\cobang{\zero}\defeq \id{\zero}\]
%\[\cobang{U \oplus P}\defeq \cobang{U}\oplus \cobang{P}\]
\begin{lemma}\label{lemma:stmat coproduct}
Let $\Cat{C}$ be a $\Cat{PCA}$-enriched category. Every object of $\stmat{\Cat{C}}$ is a natural and coherent monoid object.
\end{lemma}


\begin{lemma}\label{lemma:stmat copca}
Let $\Cat{C}$ be a $\Cat{PCA}$-enriched category. Every object of $\stmat{\Cat{C}}$ is a natural and coherent co-pca object.%, i.e., it is equipped with a co-pca structure $(P,\diagp{P},\bangp{P})$ satisfying the axioms in Figures~\ref{fig:co-pca axioms},\ref{fig:copcacoherence}.
\end{lemma}

\begin{proposition}\label{prop:cmatrixcb}
     $\stmat{\Cat{C}}$ is a convex biproduct category.
\end{proposition}
\begin{proof}
    By Lemmas~\ref{lemma:stmat coproduct} and~\ref{lemma:stmat copca}, every object of $\stmat{\Cat{C}}$ carries natural and coherent monoid and a co-pca structures. 
Thanks to Proposition~\ref{prop:A}, we can now prove the statement by simply showing that $\stmat{\Cat{C}}$ has jointly monic projections.    
    
%where $P=\bigoplus_{k=1}^n U_k$, 

Consider two objects $Q_1=\bigoplus_{k=1}^{m_1} U_k$ and $Q_2=\bigoplus_{k=1}^{m_2} V_k$. Recall that $\pi_1\colon Q_1\piu Q_2 \to Q_1$ and $\pi_2\colon Q_1\piu Q_2 \to Q_2$ are 
$(\id{Q_1}\piu\, \bangp{Q_2})$ and $(\,\bangp{Q_1}\piu \id{Q_2})$ respectively, which by \eqref{eq:matsmc} and \eqref{eq:indmoncopca} are the $(m_1\times (m_1+m_2))$ and $(m_2\times (m_1+m_2))$ matrices displayed below.
%
\[\pi_1=\begin{pNiceMatrix}
1\cdot{\id{U_1}}  	& \emptyset & \Cdots & \emptyset & \emptyset  	& \emptyset & \Cdots & \emptyset \\
\emptyset  &   & \Ddots & \Vdots & \emptyset  &   &  & \Vdots \\	
\Vdots & \Ddots &   & \emptyset & \Vdots &  &   & \emptyset \\
\emptyset  & \Cdots & \emptyset  & 1\cdot{\id{U_{m_1}}}  & \emptyset  & \Cdots & \emptyset  & \emptyset
\CodeAfter
\line{1-1}{4-4}
\line{1-5}{4-8}
\tikz \draw[thin,dotted] (2-5.east) -- (4-7.north west);
\tikz \draw[thin,dotted] (1-6.east) -- (3-8.north west);
\end{pNiceMatrix}\]
%
\[ \pi_2=\begin{pNiceMatrix}
\emptyset  	& \emptyset & \Cdots & \emptyset & 1\cdot{\id{V_1}}  	& \emptyset & \Cdots & \emptyset \\
\emptyset  &   & \Ddots & \Vdots & \emptyset  &   &   & \Vdots \\	
\Vdots & \Ddots &   & \emptyset & \Vdots &  &   & \emptyset \\
\emptyset  & \Cdots & \emptyset  & \emptyset  & \emptyset  & \Cdots & \emptyset  & 1\cdot{\id{V_{m_2}}}
\CodeAfter
\tikz \draw[thin,dotted] (2-5.south east) -- (4-7.north west);
\tikz \draw[thin,dotted] (1-6.south east) -- (3-8.north west);
\line{1-1}{4-4}
\line{1-5}{4-8}
\end{pNiceMatrix}\]
%
%which correspond to $(\id{Q_1}\piu\, \bangp{Q_2})$ and $(\,\bangp{Q_1}\piu \id{Q_2})$ respectively. 
Consider $M,N \colon \bigoplus_{k=1}^n W_k \to Q_1 \oplus Q_2$.
Both $M$ and $N$ are $(m_1+m_2)\times n$ matrices that can be viewed as block matrices $M=\begin{pmatrix} M_1 \\ M_2 \end{pmatrix}$ and $N=\begin{pmatrix} N_1 \\ N_2 \end{pmatrix}$
where $M_1,N_1$ are $(m_1\times n)$ stochastic matrices and $M_2,N_2$ are $(m_2\times n)$ stochastic matrices. If $M;\pi_1 = N;\pi_1$ and $M;\pi_2 = N;\pi_2$, by matrix multiplication it follows that $M_1=N_1$ and $M_2=N_2$ and hence $M=N$. 
\end{proof}



Just as categories of matrices are freely generated finite biproduct categories \cite{mac_lane_categories_1978}, similarly  $\stmat{\Cat{C}}$ is the convex biproduct category freely generated by $\Cat{C}$. We  illustrate this below. 


First, for all $\Cat{PCA}$-enriched functors $F\colon \Cat{C}\to \Cat{D}$, one can define $\stmat{F}\colon \stmat{\Cat{C}}\to \stmat{\Cat{D}}$ as the functor mapping an object $\bigoplus_{k=1}^n U_k$ to $\bigoplus_{k=1}^n F(U_k)$ and mapping a matrix
$M\colon\bigoplus_{k=1}^n U_k\to \bigoplus_{k=1}^m V_k$ with entries $M_{ji}= p_{ji}\cdot f_{ji}$  into the matrix with entries $\stmat{F}(M)_{ji}\defeq p_{ji}\cdot F(f_{ji})$.%=F(p_{ji}\cdot f_{ji})
\begin{proposition}\label{prop:stmatfun}
    Let $F\colon \Cat{C}\to \Cat{D}$ be a $\Cat{PCA}$-enriched functor. $\stmat{F} \colon\stmat{\Cat{C}}\to \stmat{\Cat{D}}$ is a morphism of convex biproduct categories.
\end{proposition}
Then, it is easy to check that the assignment $\Cat{C} \mapsto \stmat{\Cat{C}}$ and $F \mapsto \stmat{F}$ provides a functor $\stmat{-}:\Cat{PCACat}\to \Cat{CBCat}$ which is left adjoint to the forgetful functor $U\colon \Cat{CBCat} \to \Cat{PCACat}$.
\begin{theorem}\label{thm:matfree}
    $\stmat{-} \colon \Cat{PCACat}\to \Cat{CBCat}$ is left adjoint to $U\colon \Cat{CBCat} \to \Cat{PCACat}$.
\end{theorem}
We conclude this section with a simple observation that will be useful later in Section~\ref{sec:probbooltapes}.
\begin{proposition}\label{prop:counit fullfaithful}
    For any convex biproduct category $\Cat{C}$, the counit of the adjunction above\\ $\epsilon_\Cat{C}\colon \stmat{U(\Cat{C})}\to \Cat{C}$ is a full and faithful morphism of convex biproduct categories.
\end{proposition}


